% !TeX program = xelatex
% ==========================================================================
% Пропоузал семестрового проєкту — курс «Штучний інтелект», УКУ.
% Це ЄДИНИЙ файл, який вам потрібно редагувати. Уся верстка й брендинг
% живуть у ucuproposal.cls — не змінюйте шрифти, поля, кеглі й кольори.
%
% Компіляція: XeLaTeX + Biber (Overleaf робить це автоматично).
% Локально:   latexmk (див. README.md).
% ==========================================================================
\documentclass[ukrainian]{ucuproposal}
% Опції класу (можна комбінувати через кому):
%   ukrainian | english  — мова документа й логотипу (за замовчуванням ukrainian)
%   slogan                — додає гасло УКУ в колонтитул титулки (за замовчуванням вимкнено)

% --------------------------------------------------------------------------
% Метадані титульної сторінки
% --------------------------------------------------------------------------
\projecttitle{AI Traversability Costmap}
% \projecttagline{Одне речення про суть проєкту (опційно — можна видалити рядок)}

\teammember{Засимович Богдан}
\teammember{Ягода Микита}
\teammember{Яковенко Олена}
% \teammember[https://linkedin.com/in/...]{Прізвище Ім'я} % опційне клікабельне посилання на профіль (LinkedIn тощо)

\mentor{Олександр Косован} % якщо ментора немає — видаліть текст і залиште \mentor{}, рядок не друкується

\track{engineering} % research | engineering — не друкується на титулці, лише для внутрішньої категоризації

\repo{https://github.com/UCUSpaceRobotics/traversability-costmap.git} % опційно; якщо репозиторію ще немає — залиште порожнім, рядок не друкується

\proposaldate{2026-09-20}

\begin{document}
\maketitlepage

\section{Проблема та мотивація}
% Що саме ви розв'язуєте і чому це важливо й цікаво в контексті ШІ?
% Орієнтовно 1 абзац.
% Опишіть проблему, яку розв'язує ваш проєкт, та її значущість.
Автономна навігація на неструктурованому рельєфі стикається з обмеженнями геометричних підходів (карт зайнятості): вони не розрізняють фізичних властивостей поверхонь. Так, візуально плаский пісок може спричинити пробуксування, а безпечна трава — здаватися перешкодою. У проєкті ми розв'язуємо цю задачу, прогнозуючи прохідність місцевості (очікуване просковзування коліс і вібрації) на основі RGB-D даних. Рішення розробляється для нашого ровера, щоб виконати завдання з автономної навігації на змаганнях European Rover Challenge (ERC). З погляду ШІ проєкт цікавий застосуванням self-supervised навчання: модель тренується без ручної розмітки, зіставляючи зображення з реальними даними IMU та одометрії під час руху. Отримана адаптивна карта вартості (costmap) буде інтегрована в стек ROS 2 Nav2, що мінімізує ризик застрягання ровера та потребу у втручанні оператора.

\section{Огляд джерел та наявних рішень}
% Що вже зроблено у світі за цією темою? Спирайтесь на references.bib,
% цитуйте через \cite{}. Орієнтовно 3-5 джерел.
% Стисло опишіть наявні підходи та роботи, на які ви спираєтесь~\cite{example2023,example2022}.
Через обмеженість класичних геометричних підходів до оцінки прохідності рельєфу, фокус у робототехніці змістився на самокероване навчання (self-supervised learning). Наприклад, система BADGR~\cite{kahn2021badgr} прогнозує ймовірність зіткнень та тряски на основі послідовності зображень і навігаційних команд. Однак вона не генерує класичної просторової карти вартості (costmap), що ускладнює її інтеграцію з стеком ROS 2 Nav2, який ми використовуємо для планування шляху.

Проблему сумісності з класичними стеками навігації розв'язують архітектури на кшталт WayFASTER~\cite{gasparino2024wayfaster}, які прогнозують рівень зчеплення коліс і формують локальну карту прохідності у вигляді зверху (Bird's-Eye View, BEV). Водночас робота «How Does It Feel?» (HDIF)~\cite{castro2023hdif} довела ефективність використання показників IMU як псевдо-істинних міток для передбачення вібрацій шасі та впливу швидкості руху на прохідність.

\section{Дані}
% Які набори даних, звідки, ліцензія. Якщо збираєте самі — як саме.
% Якщо дані не потрібні — поясніть чому.
% Опишіть джерела даних, їх обсяг і ліцензію.
Основним джерелом даних для навчання моделі стане власний набір, зібраний за допомогою нашого ровера. Збір здійснюватиметься шляхом ручного керування на різних типах рельєфу. 

Датасет складатиметься із ROS bags, що містять:
\begin{itemize}
    \item RGB-D дані та Pointcloud з стереокамери;
    \item дані з IMU для розрахунку рівня вібрацій;
    \item показники одометрії для обчислення коефіцієнта просковзування.
\end{itemize}

\section{Підхід і методи}\label{sec:approach}
% Які алгоритми/архітектури плануєте використати, чи спираєтесь на
% існуючі реалізації та як плануєте їх покращити.
% Опишіть плановий підхід і методи.

Наш підхід базується на парадигмі self-supervised learning. Роботу над проєктом поділено на два етапи: створення базової (baseline) моделі та її подальше вдосконалення.

Для реалізації базової версії ми використаємо спрощений пайплайн: 
вхідні RGB-D кадри зі стереокамери трансформуватимуться у вигляд зверху (Bird's-Eye View) за допомогою алгоритму Inverse Perspective Mapping. Для обробки цих зображень буде використана базова згорткова нейромережа. Модель прогнозуватиме карту вартості (costmap) на основі очікуваного просковзування та рівня вібрацій. Навчання плануємо здійснювати за допомогою маскованої функції втрат, яка оцінюватиме помилку лише на тих ділянках карти, якими фактично проїхав ровер. Готова модель працюватиме як кастомний плагін у стеку ROS 2 Nav2.

Після успішного розгортання базової версії на ровері, систему можна буде покращити наступними методами:
\begin{itemize}
    \item \textbf{Двоканальна карта вартості:} модифікація моделі для повернення двоканальної карти вартості, де канали будуть окремо враховувати просковзування і вібрації. Це дозволить зручно поєднувати ці дані в Nav2 стеку і дасть можливість налаштовувати вплив кожного з них на загальну costmap.
    \item \textbf{Використання Pointcloud:} перехід від спрощеного алгоритму IPM до побудови локальної карти висот на основі просторової хмари точок (point cloud) зі стереокамери. Це дасть змогу мережі безпосередньо розуміти 3D-геометрію рельєфу, таку як круті нахили, перепади висот або великі камені.
    \item \textbf{Врахування динаміки:} додавання поточної швидкості ровера як додаткового входу моделі, оскільки фізичний вплив рельєфу (особливо вібрації) сильно залежить від темпу руху.
\end{itemize}

\section{Оцінювання результатів}
% Метрики, baseline, спосіб порівняння, якісні очікування.
% Опишіть метрики та baseline, з яким порівнюватиметесь.

\section{План і ризики}

\textbf{План роботи.} \\
Процес розробки поділено на два етапи.

\textbf{До попереднього захисту (Midterm):} 1) налаштування симуляції (Gazebo) та збір синтетичних даних; 2) реалізація пайплайну Inverse Perspective Mapping (BEV) та автоматичної розмітки даних; 3) створення і навчання базової нейромережі з маскованою функцією втрат; 4) тестування генерації costmap у симуляції.

\textbf{До фінального захисту:} 1) збір реального датасету та донавчання моделі; 2) тестування на фізичному ровері; 3) інтеграція costmap у ROS 2 Nav2; 4) впровадження покращень (за наявності часу), запропонованих у розділі~\autoref{sec:approach}.

\textbf{Ключові ризики та План Б:}
\begin{itemize}
    \item \textbf{Дані та сенсори:} проблеми зі збором або якістю реальних даних. \textit{План Б:} збільшити частку симуляції або використати датасет TartanDrive~\cite{triest2022tartandrive}.
    \item \textbf{Апаратна частина:} поломка ровера (двигунів, комп'ютера). \textit{План Б:} працювати в симуляції до завершення ремонту або використати іншого робота в УКУ.
    \item \textbf{Швидкодія:} брак обчислювальних потужностей бортового ПК. \textit{План Б:} зменшити роздільну здатність вхідних зображень або винести обчислення на віддалений комп'ютер через мережу ROS 2.
\end{itemize}

\section{Contribution statement}
% Хто за що відповідає. Відсотки не обов'язкові — якщо хочете їх вказати,
% додайте в квадратних дужках: \contribution[50\%]{Прізвище Ім'я}{...}
\begin{contributions}
\contribution{Засимович Богдан}{Налаштування симуляції, збір даних, інтеграція моделі в стек ROS 2 Nav2, деплой на ровер}
\contribution{...}{Створення архітектури нейромережі, вибір функції втрат, навчання та валідація моделі}
\contribution{...}{Парсинг сирих даних, перетворення зображень у вигляд зверху (IPM), автоматична розмітка датасету}

\end{contributions}

% Примітка: декларація використання AI-інструментів (AI usage) до цього
% пропоузалу не входить — вона живе в README.md вашого репозиторію.

\ucuprintbibliography

\end{document}
